Business automation and agent benchmarks increasingly rely on \emph{executable specifications}. Such a specification pairs natural-language requirements and policies, which tell an executor what to do, with executable checks that decide whether the result is right. These checks are usually treated as ground truth, yet recent audits of widely used agent benchmarks show that they are often wrong: they demand outcomes the requirements never state, accept behaviours the requirements forbid, or contradict each other. Finding such defects is hard. Static review of the text never sees the system run. Purely statistical audits cannot tell a bad specification from a weak model. Purely symbolic analysis of the checks cannot see the requirements at all.

We present \specdx, an automatic method that diagnoses executable specifications by running them with a population of local, open-weight LLM interpreters drawn from several model families. Each run is treated as one sample of how a reasonable executor reads the requirements. The core of \specdx is an \emph{anchored cell model}: a latent-class model in which the specification itself is one fallible rater, each checked value has a latent \emph{set} of admissible readings, and each interpreter has an accuracy that depends on its ability and on the difficulty of what it reads. We prove that, under stated assumptions, the defect form of each cell is identifiable from how its readings change with interpreter accuracy, and we give finite-sample bounds. \specdx combines this statistical evidence with logic. It enumerates minimal unsatisfiable subsets from the runs' violation sets and refutes candidates with witness states, localises ambiguities to sentences of the requirements by delta debugging, and chooses informative scenarios by Thompson sampling. It reports four typed defects: contradictions, hidden requirements, ambiguities and weak checks.

We evaluate \specdx on a procure-to-pay specification with seeded and real defects, under a pre-registered protocol. It is compared with seven published baselines that run on the same specifications, models, runs and scoring. [Results to be added: typed F1, real-defect recall and cost-effectiveness.]
